\documentclass[11pt]{article}
\usepackage[T1]{fontenc}
\usepackage[utf8]{inputenc}
\usepackage{lmodern}
\usepackage[margin=1in]{geometry}
\usepackage{amsmath,amssymb,amsthm,mathtools}
\usepackage{microtype}
\usepackage{booktabs,tabularx,longtable,array}
\usepackage{enumitem}
\usepackage{graphicx}
\usepackage[authoryear,round]{natbib}
\usepackage{url}
\usepackage{xcolor}
\usepackage[colorlinks=true,allcolors=blue!45!black]{hyperref}
\hypersetup{pdftitle={Dynamic Agent Regimes: Causal Evaluation, Experimental Design, and Policy Improvement},pdfauthor={DTR-AgentEvals working manuscript}}
\newcommand{\E}{\mathbb{E}}
\renewcommand{\P}{\mathbb{P}}
\DeclareMathOperator{\Var}{Var}
\DeclareMathOperator{\Cov}{Cov}
\DeclareMathOperator{\TV}{TV}
\newcommand{\ind}{\mathbf{1}}
\newcommand{\norm}[1]{\left\lVert#1\right\rVert}
\newcommand{\abs}[1]{\left\lvert#1\right\rvert}
\providecommand{\tightlist}{\setlength{\itemsep}{0pt}\setlength{\parskip}{0pt}}
\newtheorem{theorem}{Theorem}
\newtheorem{proposition}[theorem]{Proposition}
\newtheorem{corollary}[theorem]{Corollary}
\newtheorem{lemma}[theorem]{Lemma}
\theoremstyle{definition}
\newtheorem{assumption}{Assumption}
\newtheorem{definition}{Definition}
\theoremstyle{remark}
\newtheorem{remark}{Remark}
\setlist{itemsep=3pt,topsep=5pt}
\setlength{\parskip}{3pt}
\setlength{\emergencystretch}{2em}
\allowdisplaybreaks[2]
\pagestyle{plain}
\title{Dynamic Agent Regimes:\\Causal Evaluation, Experimental Design, and Policy Improvement for Model-Switching Agents}
\author{Working manuscript\\\small Author list and affiliations to be finalized}
\date{September 20, 2026}

\begin{document}
\maketitle
\begin{abstract}
Choosing a language model inside an agent changes the history on which later choices depend. Evaluation of a switching policy therefore requires an intervention on the joint trajectory distribution. We formulate routing as a dynamic regime over versioned model and harness configurations. Under sequential exchangeability, support and stable execution kernels, model-choice probabilities identify policy values without token-level likelihood ratios. We give a full-history doubly robust construction, its exact drift and inference conditions, and distinguish frozen-policy from behavior-adaptive targets. Further results connect prospective routing opportunities, selection-aware branch validation, cost-constrained allocation and execution-kernel sensitivity. These specialize established longitudinal inference, off-policy evaluation and sampling arguments; they do not guarantee useful precision or policy improvement. A descriptive open-weight coding study includes 2,640 randomized evaluation episodes and 3,960 fresh policy executions on 330 fixed tasks. Its learned router reduces to a fixed repair schedule, while limited visible feedback suppresses many repair opportunities. An offline/live discrepancy persists without an identified mismatch in the inspected assignment or estimation calculations. The branch comparison retains unresolved joint inference under the fixed-benchmark design. These findings narrow the current empirical contribution to an auditable evaluation case study. New model experiments and the remaining statistical validation are deferred.
\end{abstract}
\noindent\textbf{Keywords:} dynamic treatment regimes; agent evaluation; model routing; off-policy evaluation; sequential randomization; branch sampling.

\section{Introduction}
\label{sec:introduction}

A language-model agent may invoke different models while solving a single task. A small model may handle routine actions, a larger model may respond to a failed tool call, and a subsequent decision may return execution to the smaller model. Evaluating such a strategy requires more than measuring the accuracy of each model on a fixed collection of prompts. A model choice can change the generated action, the environment state, the information available at the next decision, and whether another decision occurs. The scientific question is therefore sequential: for a specified population of tasks and an execution environment, what would happen if the agent followed a different rule for choosing models throughout execution?

This question arises in an active routing literature. Preference-based query routing and causal learning from observational query feedback address model selection for individual requests \citep{ong2025routellm,tsiourvas2025causal}. Multi-round routing, budget-aware agent routing, and routing embedded within an agent harness already address sequential model selection \citep{zhang2025routerr1,zhang2026baar,liu2026harness}. The Replay Gap directly motivates evaluation of the downstream consequences of changing a model during execution \citep{gonuguntla2026replay}. These studies establish that routing itself, including routing within an agent, is not a new problem. They motivate a complementary question: under what conditions can previously collected agent trajectories support a statistically interpretable comparison between routing policies?

A retrospective comparison between tasks that switched models and tasks that did not generally answers a different question. Failures can trigger escalation, so the switched group may contain harder tasks or histories already affected by earlier model choices. Conditioning only on the final switching indicator does not account for this feedback. Similarly, replacing a model label while preserving the recorded downstream observations can hold fixed variables that the intervention would change. A valid evaluation must specify both the routing decisions being changed and the subsequent execution mechanisms allowed to respond. This requirement also clarifies when replay is appropriate: a restored environment followed by an actual new continuation is an experiment, while an unchanged recorded future requires a separate invariance justification.

Dynamic treatment regimes and off-policy evaluation provide the relevant mathematical foundation. Longitudinal causal inference formalizes interventions in which later treatment choices depend on variables affected by earlier treatments \citep{robins1986new,murphy2003optimal}. Sequentially randomized designs make the assignment mechanism explicit \citep{murphy2005smart,klasnja2015mrt}. History-based reinforcement-learning estimators provide closely related value-identification and doubly robust estimation results \citep{jiang2016dr,thomas2016data,kallus2020drl}. These frameworks share the essential identification problem; adopting the language of treatment regimes does not confer an ability to handle feedback that off-policy evaluation lacks. Our formulation uses both perspectives to connect the deployed intervention, the data-collection design, and the scope of an inferential claim.

We define a \emph{dynamic agent regime} as a stochastic rule for selecting a versioned model configuration at prospectively specified eligible decisions. The configuration and execution contract include the model weights, quantization, prompt and context processing, decoding, tool interface, and relevant environment behavior. For a policy $\pi$, its value is $V(\pi)=E[G^{\pi}]$, where $G^{\pi}$ is a prespecified bounded return from a genuine execution under that policy. Terminal task success is one example; a declared utility incorporating measured resource use is another. The task distribution and resource horizon are part of the estimand. A claim about one benchmark or model version does not automatically extend to another.

This intervention definition permits model selection to be treated as a macro action. When the conditional execution kernels are unchanged, the likelihood ratio between target-policy and logging-policy trajectories contains the target-to-logging routing probabilities; the intervening generation and environment factors cancel. The analysis consequently does not require reconstructing token-level likelihoods. It does require that the logged probabilities describe the decisions actually executed, including availability restrictions and fallbacks. We establish an exact reduction from the fine execution process to blocks separated by eligible decisions, retaining complete boundary histories. The reduction applies when eligibility is determined from information available before assignment and intervening execution mechanisms agree. Its statistical horizon counts eligible opportunities, including decisions that retain the current model, rather than only observed switches.

For this observation model, we give self-contained identification and inference results using established longitudinal methods. The analysis includes policy-ratio weighting, the fixed-policy influence function, the doubly robust remainder, and conditions for cross-fitted uncertainty quantification \citep{bang2005dr,kallus2020drl}. Known sequential randomization removes the need to estimate assignment probabilities but does not remove variance caused by poor overlap. We therefore consider supported stochastic changes around a reference router, while distinguishing fixed-reference policies from interventions defined by an unknown behavior distribution \citep{kennedy2019incremental}. A finite collection of policies frozen before independent evaluation admits a conservative simultaneous improvement certificate under bounded-score conditions. This certificate concerns average policy value within the stated population and supported policy class.

We then address two design issues that arise when validating agent policies. First, restored branches can provide paired continuation outcomes at selected histories, but their sampling distribution need not match the histories encountered under a target policy. We derive a target-prefix-weighted augmented branch score, its variance, and a cost-constrained allocation rule within a specified class of independently selected branch experiments. The allocation has the classical variance-and-cost structure associated with stratified sampling \citep{neyman1934representative}; optimal data collection for policy evaluation is itself an established research topic \citep{li2023allocation}. The contribution here is the explicit connection between prefix occupancy, selective continuation measurement, and the value contrast being validated. It does not identify a full policy value from arbitrary branches.

Second, a policy may be evaluated under one execution contract and deployed under another. We provide a finite-horizon sensitivity bound for bounded complete-trace returns when the initial distribution and routing policy agree and stagewise execution kernels differ by declared uniform total-variation bounds. This is a full-history adaptation of simulation-lemma and coupling arguments \citep{kearns2002near,lobel2024simulation}. It quantifies the implication of specified kernel discrepancies; it does not infer those discrepancies from limited logs or identify the value of an unsupported model update.

The paper's contribution is thus an explicitly scoped evaluation design, with mathematical results connecting
prospective routing opportunities, selection-aware branch validation and execution-contract sensitivity.
The general identification and influence-function machinery is inherited; the branch allocation and drift
arguments specialize established principles. A descriptive open-weight coding study exposes a learned fixed
repair schedule, restricted decision opportunities and an unresolved offline/live discrepancy. It does not
demonstrate history-adaptive improvement, estimator superiority or calibrated uncertainty. The theoretical
analysis, this case study and a prospective validation protocol together identify what remains necessary to
determine whether the design provides useful precision or a resource advantage over direct execution.

% Core theory, expanded from the repository foundation and audited 2026-09-19.
\section{Statistical experiment and intervention boundary}\label{statistical-experiment-and-intervention-boundary}

\subsection{Unit, timeline, and full history}\label{unit-timeline-and-full-history}

Initially suppose episodes are independent and identically distributed. An episode starts with task features and context \(H_1=X\), then generates

\[
O=(H_1,A_1,R_1,L_2,A_2,R_2,L_3,\ldots,A_T,R_T,L_{T+1}).
\]

All history and observation spaces are standard Borel, policies and kernels are measurable, and regular conditional laws exist. Action sets are finite. Whenever an optimizing action is selected, ties are resolved by a fixed measurable rule. Conditional regressions are defined on target-relevant support; arbitrary measurable versions may be assigned elsewhere because zero-occupancy contributions vanish.

Here \(H_t=(X,A_1,R_1,L_2,\ldots,A_{t-1},R_{t-1},L_t)\) is the \textbf{entire available history before routing at decision \(t\)}. Histories can contain text, tool responses, error indicators, elapsed time, cumulative cost, and previous model choices. No Markov assumption on a short state embedding is used. \(A_t\) is a finite-valued macro action: the choice of a frozen model/configuration for the next prespecified execution block. \(L_{t+1}\) contains the resulting generated output and environment observations, and \(R_t\) is the reward or utility accumulated in that block. The next history is the deterministic concatenation of these variables.

An execution block might be one model invocation, one tool-use cycle, or a prespecified subtask. Choose it before logging; changing the block definition changes the intervention. An intervention on routing does not automatically intervene on each generated token or each low-level tool action.

Let \(T<\infty\) be a fixed maximum number of decision opportunities. Early completion is an absorbing state: after the observed stopping time \(\tau\le T\), the only available action is a no-op with logging and target probability one, the state remains absorbed, and all subsequent rewards are zero. Emit a terminal success/failure reward at \(\tau\), or at \(T\) if the task reaches the cap. Thus

\[
G=\sum_{t=1}^{T}R_t
\]

can represent terminal task quality, cumulative resource use, or a prespecified utility. A cap-induced failure is part of the outcome definition. A crash that loses the final outcome is missing data and requires an additional observation model; it is not automatically absorption. An unbounded horizon is outside the theorems below.

\subsection{Macro actions are stochastic interventions on a fixed kernel}\label{macro-actions-are-stochastic-interventions-on-a-fixed-kernel}

Write the observed law as

\begin{equation}\label{eq:core1}
dP(O)=dP_1(H_1)\prod_{t=1}^{T}
b_t(A_t\mid H_t)\,dK_t(R_t,L_{t+1}\mid H_t,A_t),
\end{equation}

where \(b_t\) is the behavior/logging router. The kernel \(K_t\) integrates over the selected model's token randomness and subsequent tool/environment randomness. Schematically, for generated content \(Z_t\),

\begin{equation}\label{eq:core2}
K_t(dr,dl\mid h,a)
=\int E_t(dr,dl\mid h,a,z)\,M_{a,t}(dz\mid h).
\end{equation}

Model identity includes weight revision, quantization where applicable, tokenizer, prompt template, decoding settings, context handling, and memory serialization. The harness includes tools, permissions, stopping rules, reward measurement, and state restoration. If a model switch also changes these, define the action as that entire bundle or standardize them. A model-name string alone does not guarantee consistency.

A target regime \(\pi=(\pi_1,\ldots,\pi_T)\) replaces \(b_t\) by \(\pi_t\) while retaining \(P_1\) and \(K_t\). A deterministic DTR \(d\) is the special case \(\pi_t(a\mid h)=\mathbf{1}\{a=d_t(h)\}\). For now, \(\pi\) is fixed independently of the evaluation sample and does not vary along statistical submodels of \(P\).

\textbf{Macro-action implication.} If the intervention preserves \(K_t\), its likelihood ratio concerns model choice, not generated-token likelihoods. Token distributions of two models need not overlap with one another. When the logger sometimes chooses each model at each relevant history, its mixture already contains the trajectories generated by those models. This does not permit evaluating an unseen model, a new model revision, or a changed harness: each changes \(K_t\), rather than just \(b_t\).

\section{Identification assumptions and target value}\label{identification-assumptions-and-target-value}

The following assumptions identify a causal value, rather than merely an observed-law functional.

\begin{enumerate}
\def\labelenumi{\arabic{enumi}.}
\tightlist
\item
  \textbf{Well-defined intervention, consistency, and no interference.} Executing a routing history produces the potential observations for that same history and configuration. Separate episodes do not affect one another through shared writable memory, changing external systems, or shared resource contention unless those mechanisms are explicitly part of the design and estimand.
\item
  \textbf{Sequential exchangeability.} Conditional on the recorded pre-routing \(H_t\), the choice \(A_t\) is independent of the relevant future potential outcomes under supported continuations. Prospective routing randomization with correctly logged probabilities supplies this condition by design. Observational routing requires that all common causes of routing and future outcomes be captured. A human operator's unrecorded assessment or a hidden router feature can violate it.
\item
  \textbf{Sequential support.} For $P^{\pi}_{H_t}$-almost every history $h$, \(\pi_t(a\mid h)>0\) implies \(b_t(a\mid h)>0\). This is a population statement, not a claim that every long text history must repeat in a finite sample. Practical near-positivity remains a serious estimation problem.
\item
  \textbf{Stable task and transition law.} The evaluation target uses the same initial task distribution and conditional execution kernels as the logged experiment, or a separately justified transport model. Model updates, workload changes, judge drift, and a new task mix are not repaired merely by adjusting routing propensities.
\item
  \textbf{Outcome observation and integrability.} Rewards are observed as defined, and expectations exist. For later finite-sample results assume \(|R_t|\le r\), and for efficiency assume the influence function is square integrable.
\end{enumerate}

The target is

\begin{equation}\label{eq:core3}
v^\pi=E_{P^\pi}[G],\qquad
dP^\pi=dP_1\prod_{t=1}^{T}\pi_t\,dK_t.
\end{equation}

The g-formula originates in longitudinal causal inference, and DTRs describe the corresponding adaptive decision rules; the same value is the object of OPE. These are overlapping theoretical traditions, not competing estimands. \citep{robins1986new}, \citep{murphy2003optimal}, \citep{jiang2016dr}.

\begin{proposition}[G-computation and importance weighting]\label{core-result-1}
Define \(V_{T+1}^{\pi}=0\), and recursively

\[
Q_t^{\pi}(h,a)=E_P\{R_t+V_{t+1}^{\pi}(H_{t+1})\mid H_t=h,A_t=a\},
\] \begin{equation}\label{eq:core4}
V_t^{\pi}(h)=\sum_a\pi_t(a\mid h)Q_t^{\pi}(h,a).
\end{equation}

Then
\begin{equation}\label{eq:core5}
v^\pi=E_P\{V_1^{\pi}(H_1)\}
=E_P(W_T^{\pi}G)
=E_P\left\{\sum_{t=1}^{T}W_t^{\pi}R_t\right\},
\end{equation}

where

\begin{equation}\label{eq:core6}
W_0^{\pi}=1,\qquad
W_t^{\pi}=\prod_{s=1}^{t}
\frac{\pi_s(A_s\mid H_s)}{b_s(A_s\mid H_s)}.
\end{equation}
\end{proposition}
\begin{proof}
Consistency and exchangeability identify each intervened conditional kernel with \(K_t\); sequential support permits the change of measure. Iterated integration of \eqref{eq:core3} gives the backward recursion and its first equality in \eqref{eq:core5}. In the ratio of \eqref{eq:core3} to \eqref{eq:core1}, \(P_1\) and every \(K_t\) cancel, leaving \(W_T^{\pi}\). This proves the full-trajectory identity. For a reward measurable immediately after decision \(t\), integrating all later ratios gives conditional expectation one on target-relevant prefixes; zero-weight prefixes contribute zero. Thus \(E(W_T^{\pi}R_t)=E(W_t^{\pi}R_t)\). Summing proves \eqref{eq:core5}. Absorbed steps contribute ratio one and reward zero.
\end{proof}

For a deterministic rule, the numerator in \eqref{eq:core6} is the indicator of matching the rule. The product \(\prod_t1/b_t(A_t\mid H_t)\) by itself is \textbf{not} the weight for a specified policy; the target numerator is essential. Likewise, self-normalized or clipped importance sampling is a different estimator and does not inherit exact unbiasedness automatically.

The same factorization proves the macro-action implication in the macro-action construction above even if \(M_{a,t}\) and \(M_{a',t}\) have disjoint supports. It is sufficient that the behavior mixture chooses each needed macro action; one never divides one model's token probability by the other model's token probability.

\section{Fixed-policy influence function and estimation}\label{fixed-policy-influence-function-and-estimation}

Suppress superscript \(\pi\) in \(Q_t,V_t,W_t\). Define the uncentered score

\begin{equation}\label{eq:core7}
\phi^{\pi}(O;Q,b)
=V_1(H_1)+\sum_{t=1}^{T}W_t
\{R_t+V_{t+1}(H_{t+1})-Q_t(H_t,A_t)\}.
\end{equation}

\begin{theorem}[Efficient influence function in the full-history model]\label{core-result-2}
For a fixed target \(\pi\), in the unrestricted dominated model \eqref{eq:core1}, provided the expression is square integrable,

\begin{equation}\label{eq:core8}
D^{\pi}(O)=\phi^{\pi}(O;Q,b)-v^{\pi}
\end{equation}

is the efficient influence function (EIF). It remains efficient when the behavior probabilities are known and fixed by design, in the otherwise unrestricted model. This efficiency statement is relative to the full-history model. Additional true Markov restrictions can yield a smaller efficiency bound and different efficient estimators; using a state embedding does not establish such restrictions. \citep{kallus2020drl}.
\end{theorem}
\begin{proof}
A regular submodel has score decomposition
\[
S=S_1(H_1)+\sum_{t=1}^{T}S_{b,t}(A_t,H_t)
+\sum_{t=1}^{T}S_{K,t}(R_t,L_{t+1},H_t,A_t),
\]

with \(E S_1=0\), \(E(S_{b,t}\mid H_t)=0\), and \(E(S_{K,t}\mid H_t,A_t)=0\). Each residual in \eqref{eq:core7} has conditional mean zero given \(H_t,A_t\), so \(E D^{\pi}=0\).

Differentiating the g-formula along this submodel gives

\begin{equation}\label{eq:core9}
\dot v^{\pi}
=E[(V_1-v^{\pi})S_1]
+\sum_{t=1}^{T}E[W_t\{R_t+V_{t+1}\}S_{K,t}].
\end{equation}

There is no behavior-score derivative because \(\pi\) is fixed. In the term for a transition at \(t\), past rewards have zero covariance with its conditionally mean-zero score, and integrating future target transitions replaces future rewards by \(V_{t+1}\). Subtracting \(Q_t\) in \eqref{eq:core9} changes nothing because \(E(S_{K,t}\mid H_t,A_t)=0\).

The baseline component \(V_1-v^{\pi}\) and the weighted transition residuals form orthogonal martingale increments in the full-history filtration. Their covariances with all other score increments are zero. Hence \eqref{eq:core9} equals \(E(D^{\pi}S)\). These components lie in the corresponding baseline and transition tangent spaces, proving that \eqref{eq:core8} is the canonical gradient. Removing the behavior tangent spaces when \(b\) is known does not change this gradient, because it is already orthogonal to them.
\end{proof}

This is the established longitudinal/OPE influence-function construction, expressed for an agent-level treatment. It does not imply that a high-dimensional history regression will be accurate with a small benchmark.

\begin{theorem}[Exact drift and the scope of double robustness]\label{core-result-3}
Let \(\bar Q_t,\bar b_t\) be arbitrary fixed fitted functions with positive denominators wherever needed, let \(\bar V_t=\sum_a\pi_t\bar Q_t\), and let \(\bar W_t=\prod_{s\le t}\pi_s/\bar b_s\). Define \(\bar\phi\) by substituting these functions in \eqref{eq:core7}. Then the following exact identity holds whenever its terms are integrable:

\begin{equation}\label{eq:core10}
\begin{split}
E_P(\bar\phi)-v^{\pi}
=\sum_{t=1}^{T}E_P\Bigg[\bar W_{t-1}
\sum_a\pi_t(a\mid H_t)
\frac{\bar b_t(a\mid H_t)-b_t(a\mid H_t)}{\bar b_t(a\mid H_t)}
\{\bar Q_t(H_t,a)-Q_t^{\pi}(H_t,a)\}\Bigg].
\end{split}
\end{equation}
\end{theorem}
\begin{proof}
Put \(e_t=\bar Q_t-Q_t^{\pi}\) and \(d_t=\bar V_t-V_t^{\pi}=\sum_a\pi_te_t\), with \(d_{T+1}=0\). The true Bellman residual has conditional mean zero even when multiplied by \(\bar W_t\). Therefore

\[
E(\bar\phi)-v^{\pi}
=E(d_1)+\sum_t E\{\bar W_t(d_{t+1}-e_t(H_t,A_t))\}.
\]

Collect adjacent \(d_t\) terms to obtain

\[
\sum_t E\{\bar W_{t-1}d_t-\bar W_te_t(H_t,A_t)\}.
\]

Conditional on \(H_t\), the second term integrates over \(A_t\) under \(b_t\). Substituting \(d_t=\sum_a\pi_te_t\) gives \eqref{eq:core10}.
\end{proof}

Two familiar consequences are exact unbiasedness if every \(\bar b_t=b_t\), regardless of the fitted outcome functions, and exact unbiasedness if every \(\bar Q_t=Q_t^{\pi}\), regardless of the fitted behavior functions. Conditional unbiasedness under known randomization requires that the fitted functions be obtained independently of the scored episode, for example by sample splitting. These are algebraic population statements; estimation without splitting need not be unbiased in finite samples. \citep{bang2005dr}.

Equation \eqref{eq:core10} also vanishes for stagewise mixed patterns when, at each \(t\), either \(\bar b_t=b_t\) or \(\bar Q_t=Q_t^{\pi}\). \textbf{The latter is the true target-continuation regression, not just a correctly specified regression of an incorrectly constructed pseudo-outcome.} In ordinary backward fitting, errors in \(\bar V_{t+1}\) can propagate into \(\bar Q_t\). A naive recursive regression does not automatically converge to \(Q_t^{\pi}\) under every mixed pattern. We therefore call \eqref{eq:core7} a longitudinal DR/OPE score and do not claim that the ordinary backward-regression fitting algorithm has the stronger algorithmic sequential-double-robustness property. Such a claim requires a suitable fitting construction and proof. \citep{luedtke2017sdr}.

For a useful rate bound, define \(\mu_t(dh,a)=P(H_t\in dh)\pi_t(a\mid h)\). If \(\bar b_t\ge\epsilon>0\) on target-supported actions and \(\|\bar W_{t-1}\|_\infty\le C_{t-1}\), Cauchy--Schwarz gives

\begin{equation}\label{eq:core11}
|E(\bar\phi)-v^{\pi}|
\le\sum_{t=1}^{T}\frac{C_{t-1}}{\epsilon}
\|\bar b_t-b_t\|_{L_2(\mu_t)}
\|\bar Q_t-Q_t^{\pi}\|_{L_2(\mu_t)}.
\end{equation}

These are sufficient, deliberately strong overlap conditions. They are not needed in this form for identification. Large products of ratios can still make estimation impractical.

\subsection{A reproducible fitting algorithm}\label{a-reproducible-fitting-algorithm}

Split independent units into a fixed number of folds. For each held-out fold, fit all nuisance functions on other folds. For each fixed policy, work backward from \(\widehat V_{T+1}=0\), regress \(R_t+\widehat V_{t+1}(H_{t+1})\) on \((H_t,A_t)\), and set \(\widehat V_t=\sum_a\pi_t\widehat Q_t\). Use the \textbf{actual logged randomization probabilities} when available; do not replace them by an estimated classifier for convenience. Score each held-out episode with \eqref{eq:core7}, then average over all episodes. Model fitting, preprocessing, representation learning using outcomes, and tuning belong entirely within training data. Additional nested splitting within training can help control nuisance overfitting but is not a substitute for the convergence assumptions below.

For comparison, include the empirical mean in directly randomized target-policy arms, sequential g-computation, per-decision IPW, and the longitudinal DR score. Unsupported regression predictions are model extrapolations, not empirical identification.

\begin{theorem}[Cross-fitted asymptotic normality]\label{core-result-4}
Suppose \(T\) and the number of folds are fixed, every fold size divided by \(n\) tends to a strictly positive constant, evaluation episodes are i.i.d., \(\pi\) is fixed, \(0<\sigma_\pi^2=E[(D^{\pi})^2]<\infty\), and for every fold:

\begin{itemize}
\tightlist
\item
  the fitted score converges to \(\phi^{\pi}(O;Q,b)\) in \(L_2(P)\);
\item
  its population drift in \eqref{eq:core10} is \(o_p(n^{-1/2})\), for example by \eqref{eq:core11} and the corresponding product-rate conditions;
\item
  score integrability and variance consistency hold, for example through bounded rewards, fixed overlap constants, and suitably bounded fitted regressions.
\end{itemize}

Then
\begin{equation}\label{eq:core12}
\sqrt n(\widehat v^{\pi}-v^{\pi})
=n^{-1/2}\sum_{i=1}^{n}D^{\pi}(O_i)+o_p(1)
\ \rightsquigarrow\ N(0,\sigma_\pi^2).
\end{equation}
\end{theorem}
\begin{proof}
Within each fold, decompose the estimation error into the empirical mean of the true EIF, the empirical process of fitted-score minus true-score, and the fitted-score population drift. Conditional on training data, the empirical-process term has mean zero and variance bounded by its \(L_2(P)\) squared error divided by fold size. It is therefore \(o_p(n^{-1/2})\). A fixed finite number of folds preserves this order without requiring independence across fitted folds. The drift is negligible by assumption. The classical i.i.d. central limit theorem applies to the true EIF average.
\end{proof}

Estimate variance with the sample variance of held-out scores, and use \(\widehat v\pm1.96\widehat\sigma/\sqrt n\) for a pointwise approximate 95\% interval. With unknown propensities, both nuisance errors of order \(o_p(n^{-1/4})\) suffice for \eqref{eq:core11}, but high-dimensional agent histories do not guarantee that rate. With known propensities the drift is zero; efficiency still requires the score to converge to the EIF.

\textbf{Point consistency does not establish confidence-interval validity.} For example, suppose a behavior fit converges to a fixed incorrect \(\bar b\), while a fitted correct outcome model has error \(O_p(n^{-1/2})\). Equation \eqref{eq:core10} then generally has drift \(O_p(n^{-1/2})\), rather than \(o_p(n^{-1/2})\). The value estimate can be consistent, but the first-order randomness of fitting Q can enter its limiting distribution. A Wald interval computed solely from the evaluation-score variance omits that contribution. Cross-fitting removes evaluation-sample overfitting; it does not remove this first-order drift. Reporting all four nuisance-misspecification cells as though DR point consistency guaranteed 95\% coverage would be incorrect. To justify such inference, derive the additional nuisance influence terms, obtain a negligible drift by stronger conditions, or use a separately justified inferential method.

\textbf{Known-randomization corollary.} If \(b\) is known exactly and the cross-fitted scores converge in \(L_2(P)\) to a fixed \(\phi^{\pi}(O;\bar Q,b)\), then \eqref{eq:core10} gives zero drift even when \(\bar Q\ne Q^{\pi}\). The proof of Theorem~\ref{core-result-4} instead gives asymptotic normality with influence function \(\phi^{\pi}(O;\bar Q,b)-v^{\pi}\), provided it has a finite nonzero variance and the empirical variance is consistent. Score-based intervals are valid under these conditions, but efficiency is not guaranteed. This is a central reason to prefer prospective logging with known randomization probabilities in the first agent experiment.

\section{Task clusters, contrasts, and resource outcomes}\label{task-clusters-contrasts-and-resource-outcomes}

\subsection{Repeated seeds and branched trajectories}\label{repeated-seeds-and-branched-trajectories}

If each task is run with several seeds, branches, or models, the independent unit is usually the \textbf{task cluster}. Randomly assigning episodes from one task to both training and evaluation leaks task information and can understate uncertainty. Split by task and, where appropriate, by related task family or source repository.

For \(J\) independent and identically distributed task clusters with a fixed number \(m\) of valid root-to-terminal runs per task, assume each run has the specified behavior-law marginal and fit nuisances outside the entire held-out cluster. Define the task-weighted estimator as the average of cluster mean scores. Under the corresponding cluster-scale score-convergence, negligible-drift, and variance-consistency conditions, the cluster influence contribution is

\begin{equation}\label{eq:core13}
D_j^{\mathrm{cluster}}=m^{-1}\sum_{r=1}^{m}D^{\pi}(O_{jr}).
\end{equation}

Use the empirical variance of cluster scores divided by \(J\), not an episode-level variance divided by \(Jm\). Dependence from paired random seeds is allowed within a cluster. Selected continuation branches do not automatically have the root-episode marginal; clustering alone cannot repair that sampling bias. Their continuation estimand and sampling weights must be derived separately before a valid task-level contribution is averaged. Unequal or outcome-dependent numbers of runs require explicitly defining task versus episode weighting and rederiving the corresponding ratio or weighted estimator. A task-weighted average and an episode-weighted average need not estimate the same population.

For a fixed benchmark with repeated seeds, uncertainty over random execution seeds is conditional on that benchmark; it is not automatically uncertainty over a wider task population. If the behavior router learns from outcomes of previous episodes, the data may require martingale or adaptive-experiment theory. The i.i.d. theorem is not a justification for that setting.

\subsection{Policy contrasts and a cost vector}\label{policy-contrasts-and-a-cost-vector}

For two frozen policies evaluated on the same logs,

\begin{equation}\label{eq:core14}
\Delta=v^{\pi_1}-v^{\pi_0},\qquad
D_\Delta=D^{\pi_1}-D^{\pi_0}.
\end{equation}

Compute the variance of the paired score difference; summing two separate variances discards covariance. For \(J\) task clusters, apply the same operation to cluster means.

It is usually better to preserve a reward vector

\[
R_t^{\mathrm{vec}}=(R_t^{\mathrm{quality}},C_t^{\mathrm{tokens}},
C_t^{\mathrm{latency}},C_t^{\mathrm{money}},C_t^{\mathrm{switch}})
\]

and estimate its value vector and covariance. A prespecified scalar utility is a linear combination, such as quality minus \(\lambda^\top\)cost. State units and \(\lambda\); selecting \(\lambda\) after inspecting evaluation results introduces additional selection. Money can be near zero for local inference while latency and energy are not. Serial stage times sum to latency; parallel execution needs an explicit elapsed-time outcome rather than assuming stage-time additivity. Quantiles and tail-risk outcomes need their own influence-function arguments.

One can optimize quality subject to expected token and latency budgets using a finite candidate set and simultaneous uncertainty bounds. This is not the same as a per-task hard budget, which must be built into the policy's allowable actions and absorbing rule.

\section{Policy learning and a finite-class improvement guarantee}\label{policy-learning-and-a-finite-class-improvement-guarantee}

\subsection{Correct optimal continuation}\label{correct-optimal-continuation}

For an unrestricted supported policy class and the specified scalar utility, Bellman's recursion on full history is

\[
V_{T+1}^{*}=0,\qquad
Q_t^{*}(h,a)=E[R_t+V_{t+1}^{*}(H_{t+1})\mid H_t=h,A_t=a],
\] \begin{equation}\label{eq:core15}
V_t^{*}(h)=\max_{a\in\mathcal A_t(h)}Q_t^{*}(h,a),\qquad
d_t^{*}(h)\in\arg\max_aQ_t^{*}(h,a).
\end{equation}

The admissible set \(\mathcal A_t(h)\) must respect support and any hard resource restrictions. The proof is backward induction: the last decision maximizes its immediate expected remaining reward; conditional on the optimal continuation at later stages, maximizing \eqref{eq:core15} is optimal at the current stage. Full history suffices; a Markov state compression is optional and requires justification.

The naive \(\arg\max_aE[Y\mid H_t=h,A_t=a]\) usually evaluates \textbf{behavior-policy continuation}, not optimal continuation. For example, at decision 1, action 0 leads to terminal mean 0.6 regardless of decision 2. Action 1 leads to terminal reward 1 if decision 2 uses action 1 and 0 otherwise. If the behavior policy chooses that latter action only 10\% of the time, the observed conditional means are 0.6 versus 0.1, so the naive rule chooses 0. The optimal continuation after action 1 gives value 1, so the true optimal first choice is 1. Both decision-2 actions can have positive probability; this failure is not merely a positivity violation.

Practical policy training can use Q-learning, policy search, or existing routing algorithms. Independent evaluation is needed regardless of the training algorithm. Inference for the value of a frozen learned policy is easier than inference for the population optimum or the identity of an optimal rule near ties; no generic regularity claim about the latter is made here.

\begin{theorem}[Held-out finite-class improvement certificate]\label{core-result-5}
Condition on a training dataset that yields \(K\) candidate policies, a baseline \(\pi_0\), and fitted outcome regressions for all of them. Evaluate on \(n\) fresh i.i.d. episodes generated under \textbf{known} \(b\), and require sequential target support for every candidate and baseline. Suppose for every candidate and the baseline:

\[
|R_t|\le r,\quad
|\bar Q_t(h,a)|\le (T-t+1)r,\quad
\pi_t(a\mid h)/b_t(a\mid h)\le c_t
\]

on realized support, with deterministic finite constants. Set \(C_t=\prod_{s\le t}c_s\) and

\begin{equation}\label{eq:core16}
M=Tr+2r\sum_{t=1}^{T}C_t(T-t+1).
\end{equation}

Compute \(\bar\phi_k\) with known \(b\), and let

\begin{equation}\label{eq:core17}
\widehat\Delta_k=n^{-1}\sum_i(\bar\phi_k(O_i)-\bar\phi_0(O_i)),
\qquad
\epsilon_n=M\sqrt{\frac{8\log(2K/\alpha)}{n}}.
\end{equation}

With conditional probability at least \(1-\alpha\), simultaneously for all \(k\),

\begin{equation}\label{eq:core18}
|\widehat\Delta_k-(v^{\pi_k}-v^{\pi_0})|\le\epsilon_n.
\end{equation}

Consequently, choosing any candidate whose lower bound \(\widehat\Delta_k-\epsilon_n\) is positive, and retaining the baseline if none is positive, certifies positive expected-utility improvement whenever a switch is made, with error probability at most \(\alpha\).
\end{theorem}
\begin{proof}
The fitted score has conditional mean \(v^{\pi_k}\) by \eqref{eq:core10}, because \(b\) is known exactly. Its initial term is bounded by \(Tr\). The absolute residual at time \(t\) is at most \(r+(T-t)r+(T-t+1)r=2(T-t+1)r\), establishing \(|\bar\phi_k|\le M\). Each score difference lies in \([-2M,2M]\), an interval of width \(4M\). Hoeffding's inequality therefore gives

\[
P\{|\widehat\Delta_k-\Delta_k|>u\mid\mathrm{training}\}
\le2\exp\{-nu^2/(8M^2)\}.
\]

A union bound over \(K\) candidates and substitution of \eqref{eq:core17} prove \eqref{eq:core18}. Selection using simultaneous bounds preserves the event.
\end{proof}

The guarantee concerns expected utility under the specified task and harness distribution; it is not a safety guarantee for every task. It permits evaluation-set selection among a frozen finite class, but not inventing new candidates after looking at the evaluation set without accounting for that adaptation. With task clusters, use independent bounded cluster mean scores and replace \(n\) by the number of clusters. Multiple outcome/budget constraints require a union bound across both candidates and outcome channels. Estimated propensities introduce drift; the theorem then requires an explicit additional bias bound. Standard cross-fitting alone does not make this finite-sample statement exact.

This deliberately elementary bound is often conservative, especially with long horizons. Its purpose is to make the conditions for a valid claim transparent. Tighter empirical-Bernstein, confidence-sequence, or localized policy-class results are extensions, not established here. Safe policy improvement using DR OPE is already part of the prior literature, including \citep{jiang2016dr}.

\section{A support-aware routing intervention}\label{a-support-aware-routing-intervention}

Deterministic ``always use the large model'' rules can be nearly unsupported. An alternative is a bounded change in the odds of choosing that model where the logger already explores it. This is an application of established incremental propensity-score interventions, not a newly invented intervention family. \citep{kennedy2019incremental}.

For two models at active decision opportunities, write \(b_t(h)=P(A_t=1\mid H_t=h)\). Absorbed opportunities retain only their no-op, ratio one, and zero policy-derivative contribution. Let \(\delta_t(h)>0\) be a prespecified measurable routing-intensity function, for example \(\delta>1\) after a recoverable error and 1 elsewhere. Define

\begin{equation}\label{eq:core19}
q_t^{\delta}(1\mid h)
=\frac{\delta_t(h)b_t(h)}{1-b_t(h)+\delta_t(h)b_t(h)},
\qquad
Z_t(h)=1+\{\delta_t(h)-1\}b_t(h).
\end{equation}

This shifts odds, not probabilities by a fixed additive amount. It preserves deterministic routing where \(b=0\) or \(b=1\); therefore it cannot identify the benefit of trying an action never tried there. It identifies a different, supported question. If the intended action is literally ``switch versus continue,'' define a binary switch indicator relative to the previous model in \(H_t\), and apply the same construction to that action. Choosing model 1 and switching are not synonyms.

\begin{proposition}[Overlap control]\label{core-result-6}
For fixed \(\delta\), the realized likelihood ratios in \eqref{eq:core19} are

\begin{equation}\label{eq:core20}
\frac{q_t^{\delta}(A_t\mid H_t)}{b_t(A_t\mid H_t)}
=\frac{\delta_t(H_t)A_t+1-A_t}{Z_t(H_t)}.
\end{equation}

If \(0<\delta_{\min}\le\delta_t(h)\le\delta_{\max}<\infty\), each ratio is at most \(\max\{1,\delta_{\max},1/\delta_{\min}\}\). Structural zeros cause no realized infinite ratio. Products may still grow exponentially with \(T\).
\end{proposition}
\begin{proof}
\(Z_t\) is a convex combination of 1 and \(\delta_t\), hence lies between their minimum and maximum. Apply this to numerator \(\delta_t\) for \(A_t=1\) and 1 for \(A_t=0\).
\end{proof}

There are three different statistical targets:

\begin{enumerate}
\def\labelenumi{\arabic{enumi}.}
\tightlist
\item
  \textbf{Known randomized reference.} If \(b\) is known and fixed by design and \eqref{eq:core19} is defined from it, \(q^{\delta}\) is a fixed policy in that model. Theorems~\ref{core-result-2}--\ref{core-result-5} apply directly.
\item
  \textbf{A learned but frozen reference.} If a reference router \(b^{\mathrm{ref}}\) is learned on independent training data and frozen, the target is \(v^{q^{\delta}[b^{\mathrm{ref}}]}\), conditional on training. Use its fixed numerator and the actual evaluation logging denominator \(b\). The target need not equal \(v^{q^{\delta}[b]}\), and the ratio bound \eqref{eq:core20} is not automatic when the two references differ. Check support explicitly.
\item
  \textbf{The unknown behavior-adaptive population target.} If the target functional itself is \(\Psi_{\delta}(P)=v^{q^{\delta}[b_P]}(P)\), perturbing \(P\) perturbs the target policy. Treating that policy as fixed gives the wrong EIF in general.
\end{enumerate}

For case 3, in the unrestricted model, the corrected EIF is

\begin{equation}\label{eq:core21}
\begin{split}
D_{\delta}^{\mathrm{adaptive}}(O)
=D^{q^{\delta}}_{\mathrm{fixed}}(O)
+\sum_{t=1}^{T}W_{t-1}^{q^{\delta}}
\frac{\delta_t(H_t)}{Z_t(H_t)^2}
\{Q_t^{q^{\delta}}(H_t,1)-Q_t^{q^{\delta}}(H_t,0)\}
\{A_t-b_t(H_t)\}.
\end{split}
\end{equation}

On structural-zero/one histories, define the additional term as zero; the unobserved contrast there need not be identified. Assume bounded rewards, fixed finite \(T\), and bounded positive \(\delta\), so the displayed terms are square integrable. At nondegenerate histories the contrasts are identified on support.

\begin{proof}[Derivation of the behavior-adaptive gradient] At a fixed history, \(\partial q_t^{\delta}(1\mid h)/\partial b_t(h)=\delta_t(h)/Z_t(h)^2\), while the derivative for action 0 is its negative. The derivative of policy value with respect to a policy perturbation at \(t\) is its target-occupancy average multiplied by the target-continuation contrast \(Q_t(1)-Q_t(0)\). Relative to observed history occupancy, that average is weighted by \(W_{t-1}^{q^{\delta}}\). Finally, a behavior-score perturbation satisfies \(\dot b_t(h)=E[(A_t-b_t(h))S_{b,t}\mid H_t=h]\). Combining these three derivatives yields the additional covariance term in \eqref{eq:core21}. The fixed-policy derivative is given by Theorem~\ref{core-result-2}. The extra terms are mean-zero action-score components; adding them gives the canonical gradient by the same tangent-space argument. This is the chain-rule construction in Kennedy's general stochastic-intervention lemmas. \end{proof}

A useful check is \(\delta_t\equiv1\). Then \(q=b\) and the adaptive target equals the observed mean return. In \eqref{eq:core21}, the binary action component becomes \(Q_t(H_t,A_t)-V_t(H_t)\); it telescopes with the fixed-policy expression to \(G-E(G)\), the EIF of the observed mean. Omitting the extra term would generally fail this check.

The adaptive target is not conventionally doubly robust against arbitrary misspecification of \(b\), because \(b\) defines the target itself. This document supplies its derivative, not a completed general cross-fitted estimator theorem or implementation for unknown \(b\). The initial empirical program should prioritize the known-randomization case and a finite prespecified \(\delta\) grid. Selecting \(\delta\) continuously on evaluation outcomes needs uniform inference or a separate evaluation sample. No fixed-policy standard error should be silently reused for case 3.

\section{Why static replay and unsupported extrapolation fail}\label{why-static-replay-and-unsupported-extrapolation-fail}

\subsection{A counterexample to holding the future history fixed}\label{a-counterexample-to-holding-the-future-history-fixed}

Let \(T=2\), no baseline covariates, \(L_2=A_1\), and terminal reward \(Y=\mathbf{1}\{A_2=L_2\}\). Consider a recorded trajectory generated by \((A_1,A_2)=(0,0)\), so its recorded intermediate state is \(L_2=0\). The target policy always takes \((1,1)\). A live target execution generates \(L_2=1\) and has reward 1. A static replay that changes the actions to \((1,1)\) but keeps the recorded \(L_2=0\) assigns reward 0. It evaluates a different intervention that freezes a treatment-affected variable.

This example does not require claiming that every replay method is invalid. Restoring a valid pre-decision snapshot and \textbf{executing the new continuation} under the target kernel is a real branch experiment. Also, replay that is correct by a specially justified environment invariance can be valid. What fails is assuming unchanged downstream states without that justification.

Simply copying the recorded outcome is a different operation from recomputing the reward after freezing the state. Outcome copying returns 1 in the example above, just as the live target does, although a policy-independent copied outcome generally cannot distinguish policies. A counterexample must therefore be attached to the specified replay operation and assumptions.

\subsection{Adaptive donor matching can discard the relevant state}\label{adaptive-donor-replay}

Consider another two-decision system, with first action fixed at zero. The first execution reveals \(U\sim\operatorname{Bernoulli}(1/2)\), and visible validation always fails at this decision. The logger then chooses \(A_2\sim\operatorname{Bernoulli}(1/2)\) independently of \(U\). Terminal hidden-test success is \(Y=\ind\{A_2=U\}\); visible validation stops the episode at decision two, including when \(Y=0\). The adaptive target chooses \(A_2=U\), so \(V(\pi)=1\).

Let an unlimited ordered bank contain independent and identically distributed donors from this same task and execution law, with order fixed before observing their records. Define replay to take the first donor's observed state, request the target action there, and return the stored outcome of the earliest donor matching the requested action prefix. Both second actions have a matching donor almost surely; no fallback is needed. If the first donor's logged action equals its state, probability \(1/2\), replay keeps that donor and returns 1. Otherwise replay selects the first later donor whose action equals the \emph{old} state \(U\). Since each donor's action is independent of its own state, selection leaves the replacement state \(U'\) an independent fair bit. Thus
\[
\E(\widehat Y_{\mathrm{replay}})
=\tfrac12\cdot1+\tfrac12\cdot\tfrac12
=\tfrac34,
\qquad V(\pi)=1.
\]
Matching the action prefix has discarded the state that determined the adaptive target's action. The discrepancy occurs with stable kernels, supported randomized actions and unlimited donors at a horizon of two. It does not imply a particular bias on any empirical benchmark.

In the same system, a target with constant second action \(c\) has value \(\P(U=c)=1/2\). Replay selects the first donor with \(A_2=c\), whose state is still a fair bit, and also averages \(1/2\). This narrow positive control explains why replay failure cannot be assumed for every policy. It supplies no general validity theorem for adaptive targets, history-dependent logging probabilities, finite-bank fallback or other forms of state substitution.

\subsection{Nonidentification without support}\label{nonidentification-without-support}

Take one decision and no covariates, with \(A=0\) and \(Y=0\) always in the logs. In world I, the potential reward for action 1 is 0; in world II, it is 1. The observed laws are identical, but the value of ``always action 1'' differs. No estimator based solely on these logs identifies that value. A language-model outcome regressor, a world model, a DR correction, or a different causal vocabulary cannot resolve this observational equivalence without an assumption or new data.

\subsection{What branches can and cannot validate}\label{what-branches-can-and-cannot-validate}

A branch created at a recorded \(H_t\) estimates a continuation effect averaged over the \textbf{branch-selection distribution of histories}, with its own randomized allocation and sampling uncertainty. If histories are selected from behavior trajectories, this is generally not the full target-policy occupancy distribution. Branches can test local calibration of \(Q_t^{\pi}\), reveal downstream divergence, and add experimental support at those particular histories.

To estimate a complete policy value from branches alone, the design must also identify or reweight the target distribution of branch histories and cover histories newly reached under target execution. A finite collection of branches does not repair missing support elsewhere. Fresh root-to-terminal executions under frozen policies provide the clearest external calibration target. Restoring filesystem state, tool sessions, hidden caches, random state, and memory correctly is part of the validity of a branch; text replay alone may not reproduce the same \(H_t\).

\section{Prospective routing opportunities, branch validation, and execution drift}
\label{sec:extensions}

This section makes three design choices explicit: when routing is allowed, which prefix population a branch experiment represents, and which changes to an execution system an evaluation tolerates. The arguments specialize familiar change-of-measure, missing-data, allocation, and coupling ideas. We do not claim that stopping-time factorization, Neyman allocation, or a simulation-lemma bound is new. The contribution is the precise contract needed to apply them to model-switching agents.

\subsection{Reduction to prospectively eligible routing opportunities}
\label{sec:eligible}

An agent may generate many tokens and tool events while changing its model at only a few opportunities. Let $s=1,\ldots,N$ index a bounded fine-time trace, and let $\mathcal F_s$ contain the complete trace immediately before a possible routing decision at fine time $s$. Terminal traces are padded. A binary eligibility indicator $E_s$ is $\mathcal F_s$-measurable and determined by a common, prespecified rule. At an eligible time a router chooses a model configuration; at an ineligible time the same prespecified continuation mechanism is used under both routers. The fine-time execution kernel, conditional on the complete history and the selected configuration, is also common. Model generation, tool execution, and memory updates belong to this kernel.

For simplicity the initial model choice is eligible at $s=1$. A policy-invariant initial block before the first opportunity can instead be integrated into the initial history and reward. Let $\sigma_k$ be the $k$th eligible time, and suppose the protocol enforces at most $K$ opportunities on every feasible trace. Missing opportunities are padded at $N+1$, with an absorbing no-op. Put $\sigma_{K+1}=N+1$. At opportunity $k$, retain the complete pre-choice history $H_k$, select $A_k$, and record the complete block $B_k$ from $\sigma_k$ up to, but not including, the choice at $\sigma_{k+1}$. The block includes its duration, observations, and accumulated reward. Thus the skeleton and its blocks reconstruct the fine-time trace. The times $\sigma_k$ can depend on previous model choices through the observed history; they need not be fixed in advance.

\begin{theorem}[Exact eligible-opportunity reduction]
\label{thm:eligible}
Assume the common eligibility and execution contract above, and sequential routing support. There are policy-invariant block kernels $\mathcal K_k(dB_k\mid H_k,A_k)$ such that, for a router $r$,
\begin{equation}
 dP^r=dP_1(H_1)\prod_{k=1}^{K}r_k(A_k\mid H_k)
        \mathcal K_k(dB_k\mid H_k,A_k).
 \label{eq:blocklaw}
\end{equation}
Generating blocks from \eqref{eq:blocklaw} and reconstructing the trace has the same distribution as fine-time execution under $r$. In particular, for target $\pi$ and behavior $b$,
\begin{equation}
 \frac{dP^\pi}{dP^b}
 =W_K=\prod_{k=1}^{K}\frac{\pi_k(A_k\mid H_k)}{b_k(A_k\mid H_k)}.
 \label{eq:eligible-ratio}
\end{equation}
If $\pi_k/b_k\le c$ on support for a constant $c\ge1$, then $W_K\le c^K$ and $\E_b W_K^2\le c^K$. Consequently, for $|G|\le B$ and $n$ independent complete traces,
\begin{equation}
 \Var_b\left\{n^{-1}\sum_{i=1}^nW_{K,i}G_i\right\}
 \le B^2c^K/n.
 \label{eq:eligible-variance}
\end{equation}
\end{theorem}

\begin{proof}
Condition on the complete history and choice at $\sigma_k$. Until the next eligible time, all routing choices are fixed by the common continuation mechanism. The distribution of the intervening fine-time events is therefore determined solely by the common execution kernels and the common eligibility rule. Stopping immediately before the next choice defines a conditional block distribution $\mathcal K_k$ independent of which router will make that next choice. Finite horizons ensure that this block is well defined. Sequential conditioning yields \eqref{eq:blocklaw}. Conversely, concatenating draws from these conditional distributions and making the specified boundary choices gives the same conditional distribution at every fine-time event, hence the same joint trace law by finite induction. Because full block contents are retained, the reconstruction preserves every measurable trace payoff.

In the ratio of the two factorizations, the initial law and all block kernels cancel. Sequential support justifies the ratio and gives \eqref{eq:eligible-ratio}; padded factors are one. Its expectation is one. Since $0\le W_K\le c^K$, $W_K^2\le c^KW_K$, and therefore $\E_bW_K^2\le c^K$. Independence and the bound on $G$ give \eqref{eq:eligible-variance}.
\end{proof}

The exponent is the number of eligible randomizations, not the number of tokens, tool calls, or observed model switches. A decision to keep the current model still contributes a likelihood ratio when its probability differs between routers. An error already observed before a decision can define eligibility; eventual failure, subsequent divergence, or whether a switch ultimately helped cannot retrospectively define it. Nor may a target silently introduce extra opportunities or a different between-opportunity controller. The theorem does not justify discarding history needed by the common block kernel. It is a finite factorization argument, not an unbounded optional-stopping result. The causal and OPE foundations are inherited from \citet{robins1986new,jiang2016dr}; the opportunity structure is a prospective experimental-design choice, related to sequential and micro-randomized designs \citep{murphy2005smart,klasnja2015mrt}.

\subsection{What a live branch experiment identifies}
\label{sec:branches}

Fix one opportunity index $k$. Draw one padded pre-choice prefix $H=H_k$ from each independent root episode under a reference prefix router $b$. Let its law be $P_B$. A target prefix router $q$ induces a probability law $P_Q$, including absorbed prefixes, on this same space. Under Theorem~\ref{thm:eligible},
\begin{equation}
 w(H)=\frac{dP_Q}{dP_B}(H)
 =\prod_{j<k}\frac{q_j(A_j\mid H_j)}{b_j(A_j\mid H_j)},
 \qquad \E_Bw=1.
 \label{eq:prefixratio}
\end{equation}
The histories are complete, so the prefix product is measurable in $H$. We assume these probabilities are known from the prospective design. Selection probabilities for branches alone do not identify this occupancy ratio. A pool containing only surviving, failed, or convenient prefixes has a different law and cannot use \eqref{eq:prefixratio} unchanged. For example, conditional on a prefix-measurable event $F$ of positive probability under both laws, the ratio becomes $w(H)P_B(F)/P_Q(F)$ on $F$. The target normalization $P_Q(F)$ need not be known even when all routing probabilities are known. Keeping one padded prefix from every root avoids this additional ratio-parameter problem.

From a faithfully restored prefix, execute two frozen continuation policies $\nu_+$ and $\nu_-$. Their remaining returns are $G_+$ and $G_-$, and $D=G_+-G_-$. Coupling their random seeds is permitted if each marginal remains a valid live continuation. Define
\begin{equation}
 \mu(H)=\E(D\mid H),\qquad v(H)=\Var(D\mid H),\qquad
 \theta=\E_Q\mu(H)=\E_B\{w(H)\mu(H)\}.
 \label{eq:branch-target}
\end{equation}
For absorbed prefixes set both remaining returns to zero. If two complete policies share prefix router $q$ and differ only in these continuations, $\theta$ is their complete-policy value contrast: their common prefix reward cancels. Otherwise it is a transported continuation contrast, not the difference between two arbitrary root policies. No claim about both continuations holding simultaneously in one physical run is needed; $\mu$ is the difference of two well-defined marginal continuation means.

To reduce branch cost, draw $S\mid H\sim\operatorname{Bernoulli}\{e(H)\}$ and observe the pair only when $S=1$. Conditional on $H$, this selection is independent of the potential branch pair. The randomization occurs before running the pair. Condition throughout on an independent training sample that fixes a regression $m(H)$ and the design $e(H)$. Assume $e>0$ wherever $w>0$ and the second moments below are finite. On $\{w=0\}$ define all weighted scores, corrections, and variance integrands involving division by $e$ to be zero, including when $e=0$. Independent root episodes use independent branch randomization and execution randomness; within-pair coupling is unrestricted subject to the correct marginals.

\begin{theorem}[Selected-branch augmentation and its exact variance]
\label{thm:branch}
For $n$ reference prefixes, including those not selected for branching, define
\begin{equation}
 U_i=w(H_i)\left[m(H_i)+\frac{S_i}{e(H_i)}\{D_i-m(H_i)\}\right],
 \qquad \widehat\theta=n^{-1}\sum_iU_i,
 \label{eq:branchscore}
\end{equation}
where the second term is zero for $S_i=0$. Then, conditional on training,
\begin{equation}
 \E(\widehat\theta)=\theta,
 \label{eq:branchunbiased}
\end{equation}
and
\begin{equation}
 n\Var(\widehat\theta)
 =\Var_B(w\mu)+\E_B\left[w^2\left\{
 \frac{v}{e}+\left(\frac1e-1\right)(\mu-m)^2
 \right\}\right].
 \label{eq:branchvariance}
\end{equation}
These are exact statements for fixed $w,m,e$, not claims of semiparametric efficiency.
\end{theorem}

\begin{proof}
On $\{w>0\}$, conditional independence implies $\E\{S(D-m)/e\mid H\}=\mu-m$, so $\E(U\mid H)=w\mu$. Its conditional second moment calculation uses
\[
 \E\{(S/e)^2(D-m)^2\mid H\}
 =\{v+(\mu-m)^2\}/e.
\]
Subtracting the squared conditional mean of the correction yields
\[
 \Var(U\mid H)=w^2\{v/e+(1/e-1)(\mu-m)^2\}.
\]
On $\{w=0\}$ the weighted identities hold trivially by the zero convention. Iterated expectation proves unbiasedness, the law of total variance gives \eqref{eq:branchvariance}, and independence divides the single-prefix variance by $n$.
\end{proof}

With finite nonzero variance, the usual i.i.d. central limit theorem and the empirical variance of the $U_i$ provide an asymptotic interval for this fixed design. Pairing can reduce or increase $v$ through its covariance but does not change $\mu$. Replacing live continuations by frozen future observations changes $D$ and its estimand. Recording only selected prefixes also omits the denominator population of \eqref{eq:branchscore}. Multiple prefixes from one root require an explicitly weighted target and cluster analysis; merely treating them as independent observations does not prove \eqref{eq:branchvariance}.

\begin{corollary}[Independent target-prefix sample]
\label{cor:two-sample-branch}
Suppose an additional independent sample $H_1^Q,\ldots,H_N^Q\sim P_Q$ is available. Then
\begin{equation}
 \widehat\theta_{\mathrm{two}}
 =N^{-1}\sum_{i=1}^Nm(H_i^Q)
 +n^{-1}\sum_{j=1}^n\frac{w(H_j)S_j}{e(H_j)}\{D_j-m(H_j)\}
 \label{eq:two-sample-branch}
\end{equation}
is unbiased, and, writing $d=\mu-m$, its conditional variance is
\begin{equation}
 \frac{\Var_Q(m)}{N}
 +\frac{\Var_B(wd)}{n}
 +\frac1n\E_B\left[w^2\left\{\frac{v}{e}+
                \left(\frac1e-1\right)d^2\right\}\right].
 \label{eq:two-sample-variance}
\end{equation}
\end{corollary}

\begin{proof}
The expectation of the first sample mean is $\E_Qm$, while the expectation of the second is $\E_B\{w(\mu-m)\}=\E_Q(\mu-m)$. The samples are independent. The first variance is $\Var_Q(m)/N$. For a summand in the second mean, its conditional mean is $wd$ and its conditional variance is the final integrand in \eqref{eq:two-sample-variance}. Total variance completes the proof.
\end{proof}

This corollary is useful only when target prefixes can actually be collected or validly transported. It does not conjure a target prefix distribution from unsupported logs. Generating target prefixes also has a cost, which must be included in any comparison of designs. We make no global efficiency claim for either estimator: known prefix ratios constrain the statistical model, and more information may be available in the source trajectories.

\begin{proposition}[Cost-constrained branch selection]
\label{prop:branch-allocation}
Keep the reference-prefix distribution, prefix sample size, and frozen $m$ fixed. Let $c(H)>0$ be the expected incremental cost of executing the branch pair, with $0<\E_Bc<\infty$. Require $\ell\le e(H)\le1$ for a fixed $0<\ell<1$ and impose $\E_B\{ce\}\le\kappa$, where $\ell\E_Bc\le\kappa\le\E_Bc$. Put
\begin{equation}
 a(H)=w(H)^2\{v(H)+(\mu(H)-m(H))^2\},\qquad \E_Ba<\infty.
 \label{eq:branch-a}
\end{equation}
For any $\lambda>0$ such that the budget is attained, the design
\begin{equation}
 e_\lambda(H)=\min\left[1,\max\left\{\ell,
       \sqrt{\frac{a(H)}{\lambda c(H)}}\right\}\right]
 \label{eq:branch-allocation}
\end{equation}
minimizes both \eqref{eq:branchvariance} and, with fixed sample sizes, \eqref{eq:two-sample-variance}. If $a>0$ almost surely and the budget is strictly between its endpoints, such a $\lambda$ exists. If $a=0$ on a set, use $e=\ell$ there. Define
\[
 \kappa_{\mathrm{sat}}=\E_B[c\{\ind(a>0)+\ell\ind(a=0)\}].
\]
For $\ell\E_Bc<\kappa<\kappa_{\mathrm{sat}}$, a positive multiplier exists; for $\kappa\ge\kappa_{\mathrm{sat}}$, $e=1$ on $a>0$ and $e=\ell$ elsewhere is optimal, with any unused budget optional. At the lower budget endpoint $e=\ell$ is optimal. When $m=\mu$, the unconstrained interior allocation is proportional to $|w|\sqrt{v/c}$.
\end{proposition}

\begin{proof}
The only term depending on $e$ in either variance is $\E_B(a/e)$. For fixed $H$ and $\lambda>0$, minimize $a/e+\lambda ce$ over $[\ell,1]$. Its derivative is $-a/e^2+\lambda c$ and its second derivative is $2a/e^3\ge0$. The clipped stationary point is \eqref{eq:branch-allocation}; when $a=0$, its minimum is at $\ell$. Thus for any feasible $e$,
\[
 \E_B(a/e_\lambda)+\lambda\E_B(ce_\lambda)
 \le\E_B(a/e)+\lambda\E_B(ce).
\]
If $\E_B(ce_\lambda)=\kappa$ and $\E_B(ce)\le\kappa$, rearrangement proves optimality. By dominated convergence, $\lambda\mapsto\E_B(ce_\lambda)$ is continuous, tending to $\ell\E_Bc$ as $\lambda\to\infty$ and to $\kappa_{\mathrm{sat}}$ as $\lambda\downarrow0$. The intermediate value theorem proves existence for an interior feasible budget. Above saturation, every informative prefix already has $e=1$, which minimizes $a/e$ pointwise; spending on $a=0$ cannot lower the variance. The lower endpoint is immediate from $c>0$ and $e\ge\ell$.
\end{proof}

This is a Neyman-type allocation calculation for the specified branch-estimation problem, not a new general theory of optimal exploration \citep{neyman1934representative,li2023allocation}. The general design depends on residual second moment, not conditional variance alone: a poor frozen continuation predictor makes under-sampled histories costly even if their branch outcomes have little intrinsic variance. The optimum is oracle because $\mu$ and $v$ are unknown. An independent pilot can estimate them and freeze an approximate allocation; unbiasedness survives any such fixed design with positive selection probabilities, while oracle optimality does not follow automatically. The optimization also holds the prefix-generation design and cost fixed and does not optimize how many root tasks to generate or how to train $m$.

\subsection{Fixed-size sampling of a recorded prefix frame}
\label{sec:fixed-frame-branches}

The independent Bernoulli selection in Theorem~\ref{thm:branch} differs from choosing a fixed number of prefixes from a completed source log. For the latter design, conditioning on that log separates uncertainty from selecting and executing branches from uncertainty in the source population and its log estimate. The following result is an elementary finite-population sampling and total-variance calculation, not a new general sampling theorem.

Condition on a complete source frame $\mathcal F$, including its log and $N\ge2$ eligible prefixes. At prefix $i$, let $d_i$ be the expected contrast between two frozen continuation policies under the specified restored-state execution law. The target is the finite-frame mean $\mu_{\mathcal F}=N^{-1}\sum_{i=1}^N d_i$. Select a simple random sample $S$ of size $2\le m\le N$, without replacement and independently of fresh continuation randomness given $\mathcal F$. For each selected prefix, observe an unbiased contrast $\widehat d_i$ of conditional variance $v_i$. Assume these contrasts are independent across prefixes given $(S,\mathcal F)$, have finite second moments, and have conditional laws unchanged by selection. These execution assumptions do not follow from distinct seed labels or matching recorded restoration checks.

For example, with $r_{ia}\ge2$ independent identically distributed outcomes under continuation arm $a\in\{0,1\}$, fixed replicate counts, and independence across arms and prefixes,
\begin{equation}
 \widehat d_i=\overline Y_{i1}-\overline Y_{i0},\qquad
 v_i=\frac{\sigma_{i1}^2}{r_{i1}}+\frac{\sigma_{i0}^2}{r_{i0}},\qquad
 \widehat v_i=\frac{s_{i1}^2}{r_{i1}}+\frac{s_{i0}^2}{r_{i0}},
 \label{eq:frame-replicates}
\end{equation}
where $s_{ia}^2$ is the unbiased arm-specific sample variance. Thus $\widehat v_i$ is conditionally unbiased for $v_i$. Coupled arms require their covariance term, while shared execution shocks across prefixes require additional covariance terms.

\begin{proposition}[Conditional variance under fixed-size prefix sampling]
\label{prop:fixed-frame-branches}
Under the preceding assumptions, define
\[
 \widehat B=\frac1m\sum_{i\in S}\widehat d_i,\qquad
 f=\frac mN,\qquad
 S_d^2=\frac1{N-1}\sum_{i=1}^N(d_i-\mu_{\mathcal F})^2,\qquad
 \overline v=\frac1N\sum_{i=1}^N v_i.
\]
Then
\begin{equation}
 \E(\widehat B\mid\mathcal F)=\mu_{\mathcal F},\qquad
 \Var(\widehat B\mid\mathcal F)=\frac{1-f}{m}S_d^2+\frac{\overline v}{m}.
 \label{eq:frame-variance}
\end{equation}
Let $s_{\widehat d,S}^2=(m-1)^{-1}\sum_{i\in S}(\widehat d_i-\widehat B)^2$. If $\E(\widehat v_i\mid\mathcal F,i\in S)=v_i$, then
\begin{equation}
 \widehat V_{\mathcal F}
 =\frac{1-f}{m}s_{\widehat d,S}^2
 +\frac f m\left(\frac1m\sum_{i\in S}\widehat v_i\right)
 \label{eq:frame-variance-estimator}
\end{equation}
is unbiased for the variance in \eqref{eq:frame-variance}. Independence between $\widehat d_i$ and $\widehat v_i$ within a prefix is not required.
\end{proposition}

\begin{proof}
Given $(S,\mathcal F)$, the mean of $\widehat B$ is $m^{-1}\sum_{i\in S}d_i$, and its variance is $m^{-2}\sum_{i\in S}v_i$. The selection indicators have inclusion probabilities $m/N$ and joint inclusion probabilities $m(m-1)/\{N(N-1)\}$. Expanding the variance of the sampled mean using these probabilities gives $(1-f)S_d^2/m$. Averaging its conditional continuation variance gives $\overline v/m$. Iterated expectation and total variance therefore prove \eqref{eq:frame-variance}.

Write $s_{d,S}^2$ for the sample variance of the fixed $d_i$ in $S$. Expansion around their sample mean gives
\[
 \E(s_{\widehat d,S}^2\mid S,\mathcal F)
 =s_{d,S}^2+\frac1m\sum_{i\in S}v_i.
\]
The inclusion probabilities also imply $\E(s_{d,S}^2\mid\mathcal F)=S_d^2$. Hence $\E(s_{\widehat d,S}^2\mid\mathcal F)=S_d^2+\overline v$, while the sampled mean of $\widehat v_i$ has expectation $\overline v$. Substitution into \eqref{eq:frame-variance-estimator} proves unbiasedness.
\end{proof}

The finite-population correction applies to between-prefix heterogeneity, not to all continuation noise. At $m=N$, fresh execution noise remains, with variance $\overline v/N$; noiseless continuations recover the usual finite-frame sampling variance. Multiplying the entire observed sample variance by $1-f$ omits $f\overline v/m$. An unbiased variance estimator alone does not establish normal-interval coverage.

\paragraph{Comparison with an estimate from the same source log.}
Let $L(\mathcal F)$ be a log-based continuation contrast measurable with respect to the source frame, and set $\widehat\Delta=\widehat B-L(\mathcal F)$. Conditional on the frame,
\begin{equation}
 \E(\widehat\Delta\mid\mathcal F)=\mu_{\mathcal F}-L(\mathcal F),\qquad
 \Var(\widehat\Delta\mid\mathcal F)=\Var(\widehat B\mid\mathcal F).
 \label{eq:frame-conditional-gap}
\end{equation}
This concerns uncertainty around a frame-specific difference, whose target is not automatically zero: the realized log estimate has already been conditioned on. Under a random-frame model with finite second moments, total variance instead gives
\begin{equation}
 \Var(\widehat\Delta)
 =\E\{\Var(\widehat B\mid\mathcal F)\}
 +\Var\{\mu_{\mathcal F}-L(\mathcal F)\}.
 \label{eq:frame-total-gap}
\end{equation}
The second term contains source-task variability and its linkage to the log estimate; Proposition~\ref{prop:fixed-frame-branches} does not supply it. Random frame sizes and multiple prefixes per task belong in that source model. Fixed-size selection can still permit a valid unconditional task-population sandwich under suitable assumptions: selection dependence alone neither proves nor disproves its validity, and replicate averages already contribute to task-score variation. Full comparison therefore requires a justified sampling argument for the declared target, rather than a mechanical correction factor applied to squared task derivatives. These conditional identities do not validate an empirical branch-minus-log interval or imply a new efficiency guarantee.

\subsection{Sensitivity to changed model or harness kernels}
\label{sec:kernel-drift}

Routing OPE presumes invariant execution kernels. The following bound states a limited alternative when that contract is only approximately satisfied. Let $\mathcal K_k$ and $\widetilde{\mathcal K}_k$ be two block-kernel systems on a common trace space, with the same initial law, eligible-opportunity contract, and frozen router $\pi$. Let total variation mean $\operatorname{TV}(P,Q)=\sup_A|P(A)-Q(A)|$. Assume standard Borel spaces so conditional couplings can be chosen measurably. A common measurable payoff $G$ takes values in $[g_-,g_+]$, with range $L=g_+-g_-$. Thus a signed payoff bounded in absolute value by $B$ has range at most $2B$, not $B$.

\begin{theorem}[Bounded execution-kernel sensitivity]
\label{thm:kernel-drift}
Suppose that for every opportunity $k$ and every relevant shared complete history and action,
\begin{equation}
 \operatorname{TV}\{\mathcal K_k(\cdot\mid h,a),
                   \widetilde{\mathcal K}_k(\cdot\mid h,a)\}
 \le\varepsilon_k,\qquad 0\le\varepsilon_k\le1.
 \label{eq:kernel-tv}
\end{equation}
Then
\begin{equation}
 |v_{\mathcal K}^{\pi}-v_{\widetilde{\mathcal K}}^{\pi}|
 \le L\left\{1-\prod_{k=1}^{K}(1-\varepsilon_k)\right\}
 \le L\sum_{k=1}^{K}\varepsilon_k.
 \label{eq:kernel-bound}
\end{equation}
If the initial laws instead differ by at most $\varepsilon_0$ in total variation, the first product can include the factor $(1-\varepsilon_0)$.
\end{theorem}

\begin{proof}
Couple equal initial states identically. While histories agree, couple routing actions identically using their common $\pi$. Conditional on a shared history and action, maximally couple the two next blocks; their probability of disagreement is at most $\varepsilon_k$. Let $C_k$ be the event that the complete histories and blocks have agreed through opportunity $k$. Then $\P(C_k\mid C_{k-1})\ge1-\varepsilon_k$ whenever the conditioning event has positive probability, so $\P(C_K)\ge\prod_k(1-\varepsilon_k)$. On $C_K$ the full traces, and hence their common payoffs, agree; otherwise their payoff difference is at most $L$. Taking expectations gives the first bound. The elementary union-product inequality gives the second. If initial laws differ, first use a maximal initial coupling whose equality probability is at least $1-\varepsilon_0$ and repeat the argument.
\end{proof}

This is a finite full-history coupling bound in the classical simulation-lemma family, not an improved general simulation lemma \citep{kearns2002near,lobel2024simulation}. It requires control of kernels at the same shared histories, including generated content, tools, rewards, and restoration state. Similarity of model names, token distributions at unrelated prompts, or a small empirical divergence rate does not establish \eqref{eq:kernel-tv}. A changed judge can be covered by including its output in the block trace and retaining the same payoff extraction function. Otherwise an additional bound on the change in the payoff itself is needed.

For two policies, let $B(\pi)$ denote the first bound in \eqref{eq:kernel-bound} for that policy. Their improvement contrast obeys
\begin{equation}
 \Delta_{\widetilde{\mathcal K}}(\pi,\pi_0)
 \ge\Delta_{\mathcal K}(\pi,\pi_0)-B(\pi)-B(\pi_0).
 \label{eq:robust-improvement}
\end{equation}
This follows by applying the value bound once to each policy and subtracting. A valid lower confidence bound for the old-system contrast can therefore be reduced by these two deterministic sensitivity allowances. A positive result certifies improvement only under the assumed kernel bounds; estimating those bounds from finite agent traces requires another uncertainty argument. In particular, the theorem supplies a sensitivity analysis rather than a data-free guarantee for an updated model.

\paragraph{Scope of the extensions.}
The three results fit together without changing their targets. Prospective eligibility defines a finite routing experiment. Branch sampling validates or estimates specific continuations under an explicitly transported prefix law. Kernel sensitivity states what must additionally be controlled to carry a value comparison to a changed execution system. None removes hidden routing confounding, identifies unsupported target prefixes, or licenses using selected branches as ordinary root episodes. Their applicability is a property of the evaluation design and recorded execution state, not of the terminology used to describe the agent.

\section{Improvement claims under a changed execution system}
\label{sec:deployment-certificate}

A confidence interval for a policy value in the logging system is not, by itself, a confidence interval for its value after a checkpoint, context processor, or tool environment changes. The kernel-drift bounds above supply one way to keep these questions separate. The following elementary consequence combines an evaluation certificate with an explicit allowance for that change.

Let $v_k$ denote the value of candidate policy $k$ in the logging system, and let $\widetilde v_k$ denote its value in a deployment system. Include the baseline as $k=0$. Suppose an evaluation procedure supplies simultaneous lower bounds $L_k$ satisfying
\begin{equation}
 \P\{v_k-v_0\ge L_k\text{ for all }k=1,\ldots,K\}\ge1-\alpha.
 \label{eq:simultaneous-deployment-input}
\end{equation}
The finite-class result above is one such procedure. Alternatively, a separately justified simultaneous inference method may be used. Assume deterministic, externally justified bounds $B_k\ge0$ such that
$|\widetilde v_k-v_k|\le B_k$ for every candidate and the baseline. For a common bounded return and specified total-variation kernel allowances, the preceding coupling result provides such $B_k$; it does not estimate them from ordinary routing logs.

\begin{corollary}[Deployment-adjusted improvement certificate]
\label{cor:deployment}
Under these conditions, with probability at least $1-\alpha$, simultaneously for all candidates,
\begin{equation}
 \widetilde v_k-\widetilde v_0\ge L_k-B_k-B_0.
 \label{eq:deployment-certificate}
\end{equation}
Consequently, select a candidate only when $L_k>B_k+B_0$, and retain the baseline otherwise. On the simultaneous coverage event, every resulting selection improves expected utility in the specified deployment system.
\end{corollary}
\begin{proof}
For each candidate,
\[
 \widetilde v_k-\widetilde v_0
 =(v_k-v_0)+(\widetilde v_k-v_k)-(\widetilde v_0-v_0)
 \ge(v_k-v_0)-B_k-B_0.
\]
Apply this inequality on the event in \eqref{eq:simultaneous-deployment-input}. Because the inequality holds for every candidate on the same event, selection among those candidates does not invalidate it.
\end{proof}

This corollary is a combination of existing concentration and perturbation arguments, not a new general safety principle. Its purpose is to make the evidentiary contract explicit: statistical precision cannot compensate for an uncontrolled change in the system being evaluated. A nominal improvement smaller than the allowed kernel drift is inconclusive for deployment. If the drift allowances themselves are estimated with simultaneous confidence $1-\beta$, the same argument gives a joint statement with probability at least $1-\alpha-\beta$ by a union bound; valid construction of those allowances is an additional task.

The guarantee concerns expected utility. It does not imply improved performance on every task, absence of harmful actions, or satisfaction of a hard resource cap. A hard cap must be built into the feasible-action and termination rules; separate outcome constraints require their own simultaneous bounds. A policy that abstains or preserves the baseline is an admissible outcome of the procedure, including when no improvement can be supported.

\section{A descriptive coding-agent case study}
\label{sec:coding-case-study}

\subsection{Design, data and evidence status}
We examine an archived coding-agent experiment to assess which parts of the evaluation contract were implemented
and which scientific questions remain unresolved. The environment uses 427 sanitized MBPP tasks and 164 HumanEval
tasks \citep{austin2021program,chen2021code}. These are public function-synthesis benchmarks, not a sample of all
agent tasks. A frozen split assigns 30 tasks to the pilot, 231 to training and 330 to evaluation; the evaluation
partition contains 239 MBPP and 91 HumanEval tasks. The intended primary inference is conditional on these fixed
benchmark tasks. This section reports observed comparisons and retrospective diagnostics, not independently
validated coverage or a completed confirmatory test of the final research claims.

The two configurations are Qwen2.5-3B-Instruct and Qwen2.5-7B-Instruct, both Q4\_K\_M, with decoding temperature
0.7, top-$p$ 0.95 and a 1,024-token output cap per call. We denote them small and large. The first model call
generates a candidate function. A second or third call is eligible only after failure of visible validation;
passing that validation causes immediate submission. Hidden benchmark tests score the submitted candidate but
do not govern deployed stopping. The visible checks were generated and certified against reference solutions
before collection, with duplicate hidden-test inputs removed. Configurations, task partitions, routing draws,
visible checks and policy definitions are frozen in the repository protocol and manifests.

The randomized log contains 4,488 episodes, eight per training or evaluation task. Within each task, the initial
assignment is blocked to four small and four large calls; later eligible assignments have probability $1/2$ per
model. The evaluation log contains 2,640 episodes. A policy learned on training tasks is frozen before the
3,960 fresh target-policy executions, comprising six policies with two runs on each of the same 330 evaluation
tasks. These repeated runs are not 3,960 independent tasks. The archived live outcomes and utility identities
have been independently reconstructed; the independent review did not rerun models or hidden-test verification.

\subsection{Observed success and resource use}
We retain the original utility $U=Y-0.01N_S-0.03N_L$, where $Y$ is hidden-test success and $N_S,N_L$ count model
calls. These unitless penalties are a declared preference, not measured money or energy. Table~\ref{tab:coding-live}
reports all six evaluated policies, without selecting a favorable subset. First-failure escalation starts small
and uses large at subsequent eligible calls. Class-tailored starts small and uses large after an assertion failure,
small after other failures. Soft escalation selects large with probability $1/2$ initially and $2/3$ subsequently. The learned
policy is discussed below. All policies share the same stopping rule and maximum of three model calls.

\begin{table}[htbp]
\centering
\small
\begin{tabular}{lrrrr}
\toprule
Policy & Success & Utility & Large calls & All calls \\
\midrule
Always small & .6106 & .5956 & .0000 & 1.5000 \\
Always large & .7167 & .6777 & 1.3000 & 1.3000 \\
First-failure escalation & .6470 & .6245 & .4152 & 1.4152 \\
Class-tailored & .6515 & .6319 & .2742 & 1.4152 \\
Soft escalation ($\delta=2$) & .6758 & .6464 & .7818 & 1.3758 \\
Learned fixed repair schedule & .7091 & .6727 & 1.1545 & 1.3333 \\
\bottomrule
\end{tabular}
\caption{Observed evaluation-cohort means, 660 episodes on 330 tasks per policy. Call counts are per episode.
These are descriptive success/resource comparisons, not an optimized frontier or a demonstration of equal quality,
policy superiority or deployment savings.}
\label{tab:coding-live}
\end{table}

The learned policy has five fewer successes than always-large among 660 executions per policy, a difference of
$-0.007576$ in success and $-0.005000$ in utility. It uses 96 fewer large calls but 118 more small calls,
or 22 more total calls. These realized differences do not establish equivalence or a positive advantage.
For the two frozen policies, revaluing the observed outcomes at penalties $c_S,c_L$ gives
\[
\widehat\Delta_U(c_S,c_L)=\frac{-5+96c_L-118c_S}{660}.
\]
The empirical crossing is $c_L=0.064375$ if $c_S=0.01$ remains fixed. If both original penalties are scaled
together, the crossing multiplier is $50/17$, giving $c_L\simeq0.088235$ and $c_S\simeq0.029412$.
This retrospective sensitivity neither changes the original endpoint nor identifies what a newly trained policy
would do at different penalties. A separate training-only refit without call penalties leaves the reachable
learned schedule unchanged.

\subsection{What the learned policy and stopping rule permit}
Independent backward-regression reconstruction reproduces all 17 frozen policy entries. On reachable states,
however, the deployed rule is simply large, then small after a failure, then large after a further failure.
All 660 live paths follow this fixed schedule at eligible stages: 542 stop after large, 16 after large/small and
102 reach large/small/large. Off-path differences in the policy table do not establish history-dependent action
selection on the policy's own trajectories. Its initial key uses only benchmark identity, preventing initial
routing by task difficulty within either benchmark. An improvement over always-small would therefore not isolate
the value of adapting to the evolving failure history.

Feedback quality further restricts the available decisions. After certification, 60 of the 330 evaluation tasks
have no surviving visible checks. On those tasks, all 120 learned and 120 always-large executions stop after
the first call, although respectively 61 and 65 fail hidden tests. Across the learned cohort, 104 first-call
submissions pass visible validation but fail hidden tests, accounting for 54.2\% of its 192 final failures.
Only 118 of 660 episodes reach a first repair. These observations identify an eligibility limitation, not the
counterfactual claim that additional calls would correct those failures. A different feedback or stopping rule
would define a new harness and require a prospective evaluation; hidden verification cannot become a deployed
stopping signal. The present data are retained in full, including zero-check tasks.

\subsection{Offline/live discrepancies and their unresolved interpretation}
Table~\ref{tab:coding-calibration} retains every live policy in the offline/live comparison. Offline entries are
the archived three-fold task-split DR estimates; the class-tailored estimate and its original fitting procedure
were independently reconstructed, while the other DR entries are transcribed from the pinned analysis summaries.
This distinction matters: a common table is not evidence that every estimator was independently reimplemented.

\begin{table}[htbp]
\centering
\small
\begin{tabular}{lrrr}
\toprule
Policy & Offline DR & Fresh live & DR minus live \\
\midrule
Always small & .5946 & .6106 & $-.0160$ \\
Always large & .7114 & .7167 & $-.0052$ \\
First-failure escalation & .6306 & .6470 & $-.0164$ \\
Class-tailored & .6149 & .6515 & $-.0366$ \\
Soft escalation ($\delta=2$) & .6689 & .6758 & $-.0069$ \\
Learned fixed repair schedule & .7260 & .7091 & $+.0169$ \\
\bottomrule
\end{tabular}
\caption{Hidden-test success estimates on the same 330 evaluation tasks. Live means have execution variability;
these differences are not bias estimates against known truth. No validated-coverage claim or ranking of estimator
accuracy follows from this single dataset.}
\label{tab:coding-calibration}
\end{table}

Class-tailored IPW success is .610606, DR is .614871 and live success is .651515. Thus the DR adjustment narrows
the discrepancy with live outcomes. Independent checks reproduce the recorded policy actions, assignment draws,
probabilities and calculations; they find no mismatch in those inspected components. A first-candidate diagnostic,
before adaptive repair, already differs between initial-small log and class-tailored live executions (.565152
versus .590909). This localization is retrospective and does not identify the cause. Finite execution variation,
differences in the execution mechanism and shortcomings in inference remain distinct possibilities. Selecting a
new estimator because it agrees with these live means would not resolve the calibration question.

The prior common-five replay comparison has mean absolute discrepancies .016061 for donor replay and .018235
for DR against noisy live means. It does not demonstrate DR superiority. Four of those deterministic targets
reduce to fixed eligible-stage schedules, and donor fallback remains present. The exact adaptive replay
counterexample in Section~\ref{adaptive-donor-replay} establishes a possible failure mechanism, not its empirical
importance in this particular dataset.

\subsection{Branch target, remaining inference and implications}
The archived branch experiment samples 200 of 564 logger-reached first-failure prefixes, spanning 103 sampled
tasks, and records two fresh continuations per model per prefix: 800 continuations in total. Each assigned model
is used throughout the remaining eligible stages. The branch contrast is .120000 and the pooled log contrast
is .134654 for stay-large minus stay-small. This is a local continuation comparison on the logger's prefix
distribution, not the learned small/large continuation, class-tailored behavior, or a whole-policy value.
Recovery after an archive-loss incident further requires an observation-process justification: retained records
and reconstructed transcript bytes do not prove that lost outcomes or execution state were exchangeable.

The primary fixed-benchmark source/selection/execution inference is not yet validated. Conditioning on the
entire observed prefix frame yields the conditional result in Section~\ref{sec:fixed-frame-branches}, but changes
the conditioning set and treats the realized log comparator as fixed. The separate iid source-population result
does not automatically apply to tasks with zero eligible prefixes or zero arm contributions. Neither the existing
derivative scale nor a variance estimate alone supplies the missing coverage argument. We consequently report
the branch points without presenting the exploratory band as a validated confidence interval.

The case study supports a narrower contribution than demonstrated adaptive improvement: a prospective routing
archive that makes design and interpretation failures inspectable. Informative deployable feedback, competitive
fixed-schedule and learned-router controls, an explicit success/resource criterion, and inference for the declared
sampling design remain necessary. Future choices will be developed outside this evaluation cohort and tested
on fresh data. A null or unfavorable comparison is retained as evidence; it is not repaired by changing the
endpoint or treating noisy estimated subgroup gains as an oracle upper bound.

\section{Prospective empirical evaluation}
\label{sec:empirical-plan}

\paragraph{Status.}
Empirical estimation and performance comparisons are deferred in this draft. This section specifies the questions future experiments must answer and the corresponding reporting structure. It contains no fitted numerical effects, benchmark rankings, or unperformed-analysis claims. Existing repository development runs are retained for provenance but will not substitute for a prespecified evaluation of the final method.

\subsection{Known-truth evaluation of the statistical claims}
Finite sequential environments will vary prior-treatment-induced confounding, sample size, eligible-decision count, policy divergence, and outcome sparsity. Exact enumeration or dynamic programming will determine policy truth. Separate cells will omit a relevant historical variable, misspecify either nuisance family, alter the execution kernels, or introduce missing terminal outcomes. A factorial design should vary horizon and overlap separately; combined stress tests will be labeled accordingly.

Comparators will include direct on-policy estimation, iterated regression, per-decision importance weighting, normalized importance weighting, and the longitudinal doubly robust estimator. A deliberately naive switch/no-switch association and a fully specified static-replay rule will serve as failure controls. When the experiment concerns branch allocation, uniform branch selection and the oracle allocation will be compared with a rule estimated on independent development data. The oracle supplies a benchmark, not an implementable competitor with privileged test information.

The replay implementation will first be checked against both positive and adaptive-negative examples with known policy values, such as Section~\ref{adaptive-donor-replay}. Donor ordering, stopping, fallback and any thresholding of stochastic targets must be explicit. Failure in a specified counterexample does not require replay to perform worse on every empirical benchmark; a descriptive discrepancy from noisy live estimates is not a repeated-sampling bias estimate or a test of equal accuracy.

The primary statistical outputs will be bias, root mean squared error, confidence-interval coverage and length, false improvement declarations, and computation. Weight tails, zero-match rates, stage-specific effective sample sizes, empty nuisance-fit cells, and numerical failures will be retained. Monte Carlo uncertainty will accompany performance summaries. An interval that fails under misspecification will be reported as a failure under the relevant assumptions rather than hidden by selective reporting or undeclared truncation.

\subsection{Prospective open-weight agent experiment}
The empirical agent study will fix a model pool, prompts, decoding settings, tool interface, execution environment, stopping rules, and objective verifier. Development tasks will establish a feasible configuration and a nondegenerate outcome distribution. After that choice, all evaluation tasks will be newly drawn or held out at the task-family level. Model weights and framework software will have separately documented licenses and immutable version identifiers.

Routing opportunities will be defined prospectively. Known assignment probabilities will be persisted before generation, including availability and budget restrictions. Candidate policies will include always-small, always-large, an observable fixed-time switch, error-triggered escalation, a supported stochastic policy, and a competitive learned router. Each rule will specify its behavior when a trigger never occurs or an action becomes infeasible. A successful episode, a budget-limited failure, and an infrastructure loss of the outcome will remain distinct statuses.

Randomized logs will supply the offline estimates. Independent root-to-terminal executions under frozen target policies will supply validation estimates with their own sampling uncertainty. Branch experiments will separately evaluate continuation contrasts under an explicit reference-prefix or target-prefix distribution. Restoration fidelity will be checked against task-relevant state, and same-configuration branches will measure execution variability. Neither a branch dataset nor a repeated seed is automatically an additional independent task.

\subsection{Policy selection, resource use, and reporting}
Training, selection, and final evaluation will be separated by task or task family. The final sample size will be chosen for a prespecified precision or power target using development estimates of task-level contrast variance; it will not be justified by the number of model calls alone. Candidate classes, cost units, primary contrasts, and multiplicity handling will be frozen before confirmatory outcomes are viewed. Local token counts, elapsed time, billed cost, and energy are distinct measurements.

The final manuscript will add three empirical components: a table of statistical operating characteristics against known truth; a comparison of offline values with fresh-policy execution values; and a success--resource frontier with independent policy-improvement contrasts. Each component will report failures and uncertainty. A single live benchmark comparison cannot establish repeated-sampling coverage, and a ranking correlation cannot establish calibrated causal value estimates. If the available computation cannot support the planned precision, the study will report the attainable precision and narrow its claim.

\section{Discussion}
\label{sec:discussion}

Dynamic agent regimes provide a precise way to ask how changing model selection changes the outcome of an agent's complete execution. The central requirement is to align the intervention, the trajectory distribution, and the evaluation design. Versioned model configurations define the available actions; prospective eligibility defines when those actions can change; and sequential assignment probabilities connect the collected trajectories to supported target policies. Under the stated assumptions, the analysis reduces to longitudinal policy evaluation. The additional results clarify how that evaluation relates to selectively observed branch continuations and to bounded changes in the execution mechanism.

This formulation does not replace existing dynamic-regime or off-policy methods. The g-formula, importance ratios, fixed-policy influence function, and doubly robust construction have established foundations. Likewise, variance-based allocation and simulation-lemma arguments precede this application. The methodological value of the present development is the explicit connection of these components to decisions that can actually be deployed and outcomes that can actually be observed. Its practical value will depend on whether those contracts can be maintained and whether the resulting estimates are informative at realistic sample sizes.

Overlap remains a central limitation. Randomization can provide support for a prespecified model pool, but a small probability of a target-compatible action sequence can still produce unstable estimates. Restricting routing to fewer eligible opportunities can shorten the likelihood-ratio product, at the cost of restricting the strategies under consideration. Supported stochastic changes similarly narrow the question; they do not reveal the value of an unobserved model or guarantee a useful improvement. Neither an effective-sample-size diagnostic nor a fitted continuation model repairs missing support. High-dimensional histories also make nuisance estimation difficult, and a convenient compressed representation need not preserve the identifying information.

The branch results require their own design contract. The selected prefix must be restored faithfully, selection probabilities must be known, and the prefix weight must correspond to the target history distribution. Pairing continuations can improve precision when it induces favorable covariance, but does not make descendants independent observations. Unequal numbers of branches can change the estimand if their contributions are averaged without appropriate weighting. Sampling structure is consequential in off-policy evaluation more generally \citep{kallus2021multiple}. The derived allocation is optimal only within the stated selection class, for the specified cost model and conditional moments. Estimating these moments in a pilot creates an additional design problem; the oracle result is not a guarantee for an adaptively tuned implementation.

The execution-drift bound has a different role from identification. It propagates assumed uniform kernel discrepancies to a worst-case value discrepancy. It does not show that model updates have small discrepancies, and checking a finite set of restored histories cannot certify a uniform bound over all target-reachable histories. A coarse reward range or accumulated discrepancy may make the bound uninformative. Nevertheless, explicitly stating the required bound distinguishes a sensitivity argument from an unsupported assertion that an earlier estimate remains valid after changing prompts, checkpoints, quantization, tools, or serving behavior.

Several practical conditions also remain outside the current model. The analysis assumes independent task units, with repeated executions and branches handled within task. Shared mutable environments, caches, or resource contention can introduce interference beyond that structure. Outcomes must correspond to the declared verifier and horizon; policy-dependent missing verification requires additional observation assumptions. A fixed benchmark supports conclusions conditional on its tasks unless a task-sampling argument is supplied. Even when policy-value inference is valid, it does not identify which model is best for each individual task or justify unrestricted search over policies using the same evaluation outcomes.

The next empirical step is to compare offline estimates with independently executed, frozen policies under a common contract. Known-truth simulations should separate algebraic correctness from coverage, support failure, and nuisance misspecification. Live studies should use models capable of nondegenerate task performance, record actual routing probabilities, and distinguish task sampling from repeated-seed variability. Restored same-model controls are needed to assess branch reproducibility. Both null improvement and cases in which direct policy execution is cheaper or more precise than off-policy evaluation are scientifically informative. The present theoretical development does not establish empirical superiority or deployment savings; the planned studies are required to determine when this evaluation design is useful.


\section*{Data, software, and manuscript status}
The editable source and supporting theoretical checks are available at
\url{https://github.com/ykzeng-yale/DTR-AgentEvals}.
This draft develops the theory, reports a descriptive coding-agent case study from archived records,
and specifies the role of future validation. Case-study outputs, pinned reconstruction scripts and their
evidence boundaries are recorded in the repository; independent reconstruction did not repeat model or
hidden-verifier execution. Previously archived synthetic and pilot runs do not replace the planned final
operating-characteristic studies. No new outcome collection was performed for this manuscript revision.
The theorem statements are proved under the conditions given here and have received internal mathematical review; this does not constitute external peer review or proof-assistant certification.

\bibliographystyle{plainnat}
\begingroup
\raggedright
\bibliography{references}
\endgroup
\end{document}
